% Shared preamble for main.tex and main-fr.tex.

% Encoding and languages
\usepackage[utf8]{inputenc}
\usepackage[T1]{fontenc}
% Numeric citations in order of appearance ([1], [2], ...), processed by BibTeX.
% natbib must be loaded before babel-french.
\usepackage[numbers,sort&compress]{natbib}
% The main language must be the document language: babel-french restyles
% lists and footnotes only when French is the main language.
\usepackage[main=\TemplateMainLanguage,english,french]{babel}
\usepackage{csquotes}
\usepackage{etoolbox}

% Page and paragraph layout
\usepackage[a4paper,left=3cm,right=2.5cm,top=3cm,bottom=3cm]{geometry}
\setlength{\parindent}{0pt}
\setlength{\parskip}{0.55em}
\renewcommand{\baselinestretch}{1.08}

% Mathematics and fonts: amsmath must precede newtxmath.
\usepackage{amsmath}
\usepackage{mathtools}
\usepackage{amsthm}
\usepackage{libertine}
\usepackage[libertine]{newtxmath}

% Figures, tables, algorithms
\usepackage{graphicx}
\usepackage{booktabs}
\usepackage{array}
\usepackage{tabularx}
\usepackage{longtable}
\usepackage[labelsep=endash]{caption}
\usepackage{subcaption}
\usepackage{pdflscape}
\usepackage{tikz}
\usetikzlibrary{arrows.meta,positioning}
\usepackage{tikz-cd}
\usepackage{mathpartir}
\usepackage{algorithm}
\usepackage{algpseudocode}
\usepackage{listings}

% Typographic polish
\usepackage{microtype}
\usepackage{xcolor}
\usepackage{fancyhdr}
\usepackage{siunitx}
\usepackage{acro}

\definecolor{LMFINavy}{HTML}{304B8A}
\definecolor{LMFITeal}{HTML}{0F766E}

% Portable code examples: listings needs no shell escape or external program.
\lstdefinestyle{TemplateCode}{
  basicstyle=\ttfamily\small,
  keywordstyle=\color{LMFINavy}\bfseries,
  commentstyle=\color{LMFITeal},
  numbers=left,
  numberstyle=\tiny\color{gray},
  numbersep=8pt,
  frame=lines,
  rulecolor=\color{LMFINavy},
  breaklines=true,
  showstringspaces=false,
  tabsize=2,
  captionpos=b
}

% Localised strings and selected metadata
\ifdefstring{\TemplateMainLanguage}{french}{%
  % babel-french normally activates :;!? to insert spaces. Active colons also
  % enter label keys such as ch:background and break cleveref. French source in
  % this template writes the required spaces explicitly, so keep punctuation
  % inactive and labels portable.
  \frenchsetup{AutoSpacePunctuation=false}
  \newcommand{\TemplateTitle}{\ThesisTitleFR}
  \newcommand{\TemplateDate}{\ThesisDateFR}
  \newcommand{\TemplateKeywords}{\ThesisKeywordsFR}
  \newcommand{\TemplatePdfLang}{fr-FR}
  \newcommand{\TemplateThesisType}{Mémoire de master}
  \newcommand{\TemplateSupervisionLabel}{Sous la direction de}
  \newcommand{\TemplateLaboratoryLabel}{Laboratoire d'accueil}
  \newcommand{\TemplateAcademicYearLabel}{Année universitaire}
  \newcommand{\TemplateLogoLabel}{Logo institutionnel facultatif}
  \newcommand{\TemplateReferentLabel}{Enseignant référent}
  \newcommand{\TemplateReferencesTitle}{Références}
  \bibliographystyle{unsrtnat-fr}
  \newcommand{\theoremname}{Théorème}
  \newcommand{\propositionname}{Proposition}
  \newcommand{\lemmaname}{Lemme}
  \newcommand{\corollaryname}{Corollaire}
  \newcommand{\assumptionname}{Hypothèse}
  \newcommand{\definitionname}{Définition}
  \newcommand{\examplename}{Exemple}
  \newcommand{\remarkname}{Remarque}
  \floatname{algorithm}{Algorithme}
  % babel-french names table captions "Table"; French theses write "Tableau".
  \addto\captionsfrench{\renewcommand{\tablename}{Tableau}}
  \renewcommand{\lstlistingname}{Extrait de code}
  \renewcommand{\lstlistlistingname}{Liste des extraits de code}
  \algrenewcommand\algorithmicrequire{\textbf{Entrée :}}
  \algrenewcommand\algorithmicensure{\textbf{Sortie :}}
  \algrenewcommand\algorithmicreturn{\textbf{retourner}}
  \algrenewcommand\algorithmicfor{\textbf{pour}}
  \algrenewcommand\algorithmicdo{\textbf{faire}}
  \algrenewcommand\algorithmicend{\textbf{fin}}
}{%
  \newcommand{\TemplateTitle}{\ThesisTitleEN}
  \newcommand{\TemplateDate}{\ThesisDateEN}
  \newcommand{\TemplateKeywords}{\ThesisKeywordsEN}
  \newcommand{\TemplatePdfLang}{en-GB}
  \newcommand{\TemplateThesisType}{Master's thesis}
  \newcommand{\TemplateSupervisionLabel}{Supervised by}
  \newcommand{\TemplateLaboratoryLabel}{Host laboratory}
  \newcommand{\TemplateAcademicYearLabel}{Academic year}
  \newcommand{\TemplateLogoLabel}{Optional institutional logo}
  \newcommand{\TemplateReferentLabel}{Academic referent}
  \newcommand{\TemplateReferencesTitle}{References}
  \bibliographystyle{unsrtnat}
  \newcommand{\theoremname}{Theorem}
  \newcommand{\propositionname}{Proposition}
  \newcommand{\lemmaname}{Lemma}
  \newcommand{\corollaryname}{Corollary}
  \newcommand{\assumptionname}{Assumption}
  \newcommand{\definitionname}{Definition}
  \newcommand{\examplename}{Example}
  \newcommand{\remarkname}{Remark}
  \floatname{algorithm}{Algorithm}
  \renewcommand{\lstlistingname}{Code listing}
  \renewcommand{\lstlistlistingname}{List of code listings}
  \algrenewcommand\algorithmicrequire{\textbf{Input:}}
  \algrenewcommand\algorithmicensure{\textbf{Output:}}
  \algrenewcommand\algorithmicreturn{\textbf{return}}
}

% Labels follow the language in force, so an abstract in the secondary language
% is labelled in that language. The French keyword label carries its own
% non-breaking space before the colon.
\newcommand{\TemplateKeywordsLabel}{\iflanguage{french}{Mots-clés~:}{Keywords:}}

% Float placement: pack floats from the top of a float page and let a page hold
% more float and less forced text, which avoids loose float pages.
\makeatletter
\setlength{\@fptop}{0pt}
\setlength{\@fpsep}{14pt}
\setlength{\@fpbot}{0pt plus 1fil}
\makeatother
\renewcommand{\topfraction}{0.95}
\renewcommand{\bottomfraction}{0.95}
\renewcommand{\textfraction}{0.06}
\renewcommand{\floatpagefraction}{0.80}
\setlength{\textfloatsep}{16pt plus 3pt minus 3pt}
\setlength{\floatsep}{14pt plus 3pt minus 3pt}
% A little extra stretch for lines that carry long inline formulas.
\setlength{\emergencystretch}{3em}

% Headers: chapter title on the right, page number centred in the footer.
\pagestyle{fancy}
\setlength{\headheight}{14.5pt}
\fancyhf{}
\fancyhead[R]{\small\nouppercase{\leftmark}}
\fancyfoot[C]{\thepage}
\renewcommand{\headrulewidth}{0.4pt}
\renewcommand{\footrulewidth}{0pt}
\fancypagestyle{plain}{%
  \fancyhf{}
  \fancyfoot[C]{\thepage}
  \renewcommand{\headrulewidth}{0pt}}

% The bibliography is an unnumbered chapter listed in the table of contents.
\renewcommand{\bibsection}{%
  \chapter*{\TemplateReferencesTitle}%
  \phantomsection\addcontentsline{toc}{chapter}{\TemplateReferencesTitle}%
  \markboth{\TemplateReferencesTitle}{\TemplateReferencesTitle}}

% Links are clickable but not coloured, so the PDF prints cleanly.
\usepackage[hidelinks,unicode=true]{hyperref}
\usepackage[capitalise,noabbrev]{cleveref}

% Theorem-like environments are defined after cleveref is loaded: cleveref must
% see them in order to name shared-counter environments correctly.
% Theorem environments share one counter, numbered within chapters
% (Theorem 3.2); equations follow the same scheme, e.g. (3.1).
\theoremstyle{plain}
\newtheorem{theorem}{\theoremname}[chapter]
\newtheorem{proposition}[theorem]{\propositionname}
\newtheorem{lemma}[theorem]{\lemmaname}
\newtheorem{corollary}[theorem]{\corollaryname}

\theoremstyle{definition}
\newtheorem{definition}[theorem]{\definitionname}
\newtheorem{assumption}[theorem]{\assumptionname}
\newtheorem{example}[theorem]{\examplename}

\theoremstyle{remark}
\newtheorem{remark}[theorem]{\remarkname}

\numberwithin{equation}{chapter}

\ifdefstring{\TemplateMainLanguage}{french}{%
  \crefname{chapter}{chapitre}{chapitres}
  \crefname{section}{section}{sections}
  \crefname{figure}{figure}{figures}
  \crefname{table}{tableau}{tableaux}
  \crefname{algorithm}{algorithme}{algorithmes}
  \crefname{theorem}{théorème}{théorèmes}
  \crefname{lemma}{lemme}{lemmes}
  \crefname{corollary}{corollaire}{corollaires}
  \crefname{assumption}{hypothèse}{hypothèses}
  \crefname{example}{exemple}{exemples}
  \crefname{remark}{remarque}{remarques}
  \crefname{proposition}{proposition}{propositions}
  \crefname{definition}{définition}{définitions}
  \crefname{equation}{équation}{équations}
  \crefname{listing}{extrait de code}{extraits de code}
  \Crefname{listing}{Extrait de code}{Extraits de code}
  \providecommand{\crefpairconjunction}{}
  \providecommand{\creflastconjunction}{}
  \providecommand{\crefpairgroupconjunction}{}
  \providecommand{\creflastgroupconjunction}{}
  \renewcommand{\crefpairconjunction}{ et }
  \renewcommand{\creflastconjunction}{ et }
  \renewcommand{\crefpairgroupconjunction}{ et }
  \renewcommand{\creflastgroupconjunction}{ et }
}{%
  \crefname{chapter}{chapter}{chapters}
  \crefname{section}{section}{sections}
  \crefname{figure}{figure}{figures}
  \crefname{table}{table}{tables}
  \crefname{algorithm}{algorithm}{algorithms}
  \crefname{theorem}{theorem}{theorems}
  \crefname{lemma}{lemma}{lemmas}
  \crefname{corollary}{corollary}{corollaries}
  \crefname{assumption}{assumption}{assumptions}
  \crefname{example}{example}{examples}
  \crefname{remark}{remark}{remarks}
  \crefname{proposition}{proposition}{propositions}
  \crefname{definition}{definition}{definitions}
  \crefname{equation}{equation}{equations}
  \crefname{listing}{code listing}{code listings}
  \Crefname{listing}{Code listing}{Code listings}
}

% Number formatting follows the language of the thesis (1.5 or 1,5).
\ifdefstring{\TemplateMainLanguage}{french}{\sisetup{locale=FR}}{\sisetup{locale=UK}}

% Project macros and abbreviations (edit macros.tex, not this file).
% Project macros and abbreviations, loaded by preamble.tex.
% Keep notation that you reuse here; delete anything the thesis does not use.

% Sets and operators
\newcommand{\N}{\mathbb{N}}
\DeclareMathOperator{\Reach}{Reach}

% Semantic brackets. These two macros draw the usual double brackets without an
% extra font; stmaryrd provides true glyphs if you prefer them. \sem{\varphi}
% sets the argument, and \sem*{\varphi} resizes the brackets to fit it.
\providecommand{\llbracket}{\mathopen{[\mkern-3.3mu[}}
\providecommand{\rrbracket}{\mathclose{]\mkern-3.3mu]}}
\DeclarePairedDelimiter{\sem}{\llbracket}{\rrbracket}

% Abbreviations used by the acro package (\ac, \acs, \acf, \printacronyms).
% Give the long form in the language of the thesis.
\ifdefstring{\TemplateMainLanguage}{french}{%
  \DeclareAcronym{dag}{short = DAG, long = graphe orienté acyclique}
  \DeclareAcronym{scc}{short = CFC, long = composante fortement connexe}
}{%
  \DeclareAcronym{dag}{short = DAG, long = directed acyclic graph}
  \DeclareAcronym{scc}{short = SCC, long = strongly connected component}
}


\hypersetup{
  pdftitle={\TemplateTitle},
  pdfauthor={\ThesisAuthor},
  pdfsubject={\ThesisProgramme},
  pdfkeywords={\TemplateKeywords},
  pdflang={\TemplatePdfLang}
}

\newcommand{\TemplateActivateLanguage}{%
  \selectlanguage{\TemplateMainLanguage}%
  \ifdefstring{\TemplateMainLanguage}{french}{\shorthandoff{:}}{}%
}
\AtBeginDocument{\TemplateActivateLanguage}
